\chapter{Обзор литературы}
\label{cha:Obzor}

Трёхмерная задача упаковки изучается на протяжении нескольких десятилетий. За это время накоплен обширный арсенал методов, который целесообразно разделить на четыре направления: конструктивные эвристики, метаэвристические подходы, точные и гибридные методы, а также методы на основе машинного обучения.

\section{Конструктивные эвристики}

Конструктивные эвристики формируют решение за один проход, последовательно размещая объекты по определённому правилу. Их главное достоинство~--- скорость, а основной недостаток~--- невозможность пересмотра ранее принятых решений.

У Karabulut и İnceoğlu~\cite{karabulut} правило DBLF сводится к выбору для очередного предмета самой <<низкой>> допустимой ячейки, а при равенстве по высоте~--- позиции, смещённой к левому краю. Эту схему они встраивают в гибридный ГА как декодер хромосомы; среди четырёх вариантов целевой функции наилучшую плотность дал критерий минимума максимальной высоты загрузки.

Wei et al.~\cite{Wei2011} описали эвристику Skyline для двумерной задачи strip-packing: свободное пространство моделируется профилем <<линии неба>>, а объект размещается в первую подходящую впадину. Метод обеспечивает сложность $O(n\log n)$ и впоследствии был обобщён на трёхмерный случай.

В работе Parreño et al.~\cite{Parreno2008} свободная часть контейнера описывается набором наложенных друг на друга max-space элементов, пересчитываемых после каждого нового блока (Maximal-Space). На бенчмарке Bischoff--Ratcliff авторам удалось побить ранее опубликованные результаты в 25 случаях из 30.

Классическая стратегия First Fit Decreasing (FFD)~\cite{CoffmanFFD1996}, при которой объекты предварительно сортируются по убыванию размера, широко применяется как базовый строительный блок более сложных алгоритмов. Практическое руководство Jylänki~\cite{Jylanki2010} систематизирует свыше двадцати двумерных эвристик размещения и предоставляет открытую C++-библиотеку RectangleBinPack.

\section{Метаэвристические подходы}

Метаэвристики итеративно улучшают одно или несколько решений, не гарантируя оптимальность, но обеспечивая высокое качество за ограниченное время.

\section{Генетические алгоритмы}

Тимофеева, Чернышева, Волков и Корелин~\cite{korelin} адаптировали генетический алгоритм для трёхмерной упаковки блоков в контейнер. Авторы применили турнирную и рулеточную селекцию, модифицированный кроссинговер и мутацию перемещением гена в начало приоритетного списка.

Zhang et al.~\cite{Zhang2024} дополнили классический ГА модулем GAN: генератор подсказывает начальные хромосомы с более удачным распределением генов, что ускоряет выход к хорошим решениям. В отчётах авторов запас пустого объёма уменьшился на 2--3~п.п. по сравнению с базовым генетическим поиском.

В~\cite{TR-MFTI-khelvas2019-boxes, khel-EnT2019} Гиля-Зетинов, Хельвас и Панкратов сопоставляют слоистые эвристики и ГА для паллетизации. Качество оценивается через долю занятого объёма и показатель устойчивости; эксперименты показывают, что итоговая укладка мало меняется при небольших погрешностях в габаритах коробок.

\section{Поиск в большой окрестности (LNS)}

Pisinger и Ropke~\cite{Pisinger2005} предложили Large Neighborhood Search (LNS)~--- метаэвристику, которая на каждой итерации разрушает часть текущего решения и восстанавливает его заново. На задачах маршрутизации и упаковки LNS продемонстрировал существенное превосходство над классическими локальными поисками благодаря способности покидать локальные оптимумы.

Ropke и Pisinger~\cite{Ropke2006} развили идею адаптивного варианта ALNS, в котором веса операторов разрушения и восстановления подстраиваются в ходе поиска. На задачах pickup-and-delivery ALNS улучшил лучшие известные решения для 40\% инстансов.

Shaw~\cite{Shaw1998} впервые сформулировал constraint-based подход к LNS: разрушаемые переменные выбираются на основе их связанности (relatedness) с <<неудачными>> частями решения. Эта идея легла в основу множества современных операторов разрушения.

Обзор Pisinger и Ropke~\cite{Pisinger2010} систематизирует архитектуру LNS, выделяя выбор размера окрестности, стратегии принятия решений (simulated annealing, threshold accepting, record-to-record) и гибридизацию с точными методами.

Ekici~\cite{Ekici2023} применил LNS к задаче bin packing с конфликтами между предметами, показав, что адаптивные операторы разрушения обеспечивают преимущество на инстансах средней и большой размерности.

\section{Другие локальные методы}

В схеме Guided Local Search Faroe, Pisinger и Zachariasen~\cite{Faroe2003} величины штрафов за превышение по осям со временем ослабевают, что помогает исследовать окрестности, близкие к оценке снизу.

Calzavara et al.~\cite{Calzavara2021} формулируют паллетизацию как задачу целочисленного программирования с ограничениями на свесы и устойчивость и решают её связкой GRASP с локальной доработкой. На промышленных выборках высота штабеля оказалась на 6\% ниже, чем у действовавших ранее регламентов укладки.

\section{Точные методы и классификация задач}

Martello, Pisinger и Vigo~\cite{Martello2000} строят ветвление с отсечениями и подмешивают эвристику L2; для конфигураций до 100 коробок метод доказывает оптимальность найденного плана. Их набор тестов до сих пор служит опорной базой для сравнений.

Queiroz et al.~\cite{Queiroz2012} разработали динамические и колонковые алгоритмы для трёхмерных guillotine-задач (knapsack, cutting-stock, strip-packing).

Amossen и Pisinger~\cite{Amossen2010} распространяют guillotine-условия на многомерный случай и сводят постановку к CP-модели: при жёстком требовании <<режущих>> плоскостей недоиспользование объёма растёт всего на 0,5\%.

Wäscher, Haußner и Schumann~\cite{Wascher2007} перерабатывают схему Dyckhoff, выделяя четыре уровня детализации для cutting \& packing; эта сетка терминов сегодня часто цитируется в новых работах.

Bortfeldt и Wäscher~\cite{Bortfeldt2013} дают обзор ограничений при контейнерной загрузке за интервал 1995--2013~гг. Отдельно Bischoff и Ratcliff~\cite{Bischoff1995} связывают выбор модели или эвристики с практикой склада: запрет смешивания грузов, локальные перегрузки по опоре и т.п.

\section{Методы машинного обучения}

Fang, Rao и Shi~\cite{Fang2023} обучают пару <<актёр--критик>> размещать полосы на ленте без размеченных примеров; цель~--- укоротить занятую длину. На их данных прирост к лучшим классическим эвристикам составил около 3~п.п., что подтверждает применимость RL к геометрической упаковке.

Zhao et al.~\cite{Zhao2021} обучают политику RL для пошаговой 3D-укладки в online-режиме; проверка устойчивости опирается на stacking-tree. В сравнении с эвристическим baseline авторы сообщают о сокращении числа контейнеров до 15\%.

В обзоре Ali et al.~\cite{Ali2022} методы 3D-упаковки делятся на пакетную (offline) и потоковую (online) обработку; отдельно разобраны ограничения постановок и показатели, связанные с задержкой решения.

\section{Смежные задачи}

Ananno и Ribeiro~\cite{Ananno2024} решали многоконтейнерную 3D-BPP с реальными ограничениями, применив двухэтапный подход (зонная декомпозиция и шаблонный поиск), сократив число контейнеров в среднем на один.

Alonso et al.~\cite{Alonso2016} объединили задачу паллетизации с задачей погрузки автомобилей в междеповской логистике, получив экономию транспортных затрат 4--7\%.

Yang et al.~\cite{Yang2023} подбирают габариты коробки для e-commerce через MIP; в расчётах авторов материал на единицу отправления сократился на 8\%, а доля пустот внутри упаковки~--- примерно на 12\%.

McIlroy~\cite{TimSort_teor} исследовал алгоритм оптимистической сортировки, эффективный для частично упорядоченных данных; подобная предварительная сортировка часто применяется в задачах упаковки при построении начальных решений.

\section{Патентный обзор}

Baumann~\cite{USPatent8423178B2} запатентовал метод автоматического формирования сборной паллеты из коробок, находящихся в нескольких однородных паллетах. Pankratov, Ehrenberg и Sweet~\cite{US11613434B2} описали систему послойного формирования паллеты с программным управлением.
